% =============================================================================
% High-Performance Computing Implementation
% =============================================================================
\label{chap:hpc}

The computational demands of processing large-scale seismic datasets 
necessitate the use of High-Performance Computing (HPC) infrastructure. This 
chapter describes the HPC implementation developed for the AI4Seism seismic 
data processing pipeline, deployed on the Leonardo supercomputer at CINECA 
(Italy).

\section{Computing Infrastructure}

The pipeline was deployed on the Leonardo supercomputer at CINECA, one of 
Europe's leading HPC facilities, under the research project IscrC\_AI4Seism\_C.

\subsection{Leonardo Cluster Specifications}
The Leonardo cluster runs Red Hat Enterprise Linux 8.7 (Ootpa) and exposes one 
GPU-based partition called Booster and one CPU-based partition called DCGP 
(Data Centric General purpose module). In this thesis work we used only the 
Booster partition.\\

The Booster partition is equipped with:
\begin{itemize}
  \item \textbf{Booster module (Atos BullSequana X2135 "Da Vinci" Blade)}:
  3,456 compute nodes, each equipped with:
  \begin{itemize}
    \item \textbf{Processor}: 1$\times$ Intel Xeon Platinum 8358 (Ice Lake), 32 
    cores at 2.60 GHz
    \item \textbf{Memory}: 512 GB DDR4 RAM
    \item \textbf{GPU Accelerators}: 4$\times$ NVIDIA Ampere A100 (64 GB HBM2e)
  \end{itemize}

  \item \textbf{Internal Network}: 200G HDR InfiniBand with Dragonfly+ topology

  \item \textbf{Storage}: Lustre parallel file system optimized for 
  high-throughput I/O operations
\end{itemize}

\textbf{Workload Manager}: SLURM 23.11 is used for job scheduling and resource 
management on Leonardo.

\subsection{Resource Management}
The cluster is managed by the SLURM (Simple Linux Utility for Resource 
Management) workload scheduler, which provides:

\begin{itemize}
  \item Job queuing and prioritization
  \item Resource allocation and accounting
  \item User-friendly interface for managing computational resources
  \item Support for both interactive and batch job execution
\end{itemize}

\section{Pipeline Architecture}

The seismic processing workflow was automated through a Makefile-based 
orchestration system, enabling reproducible execution of the complete 
processing chain.

\subsection{Processing Stages}
The pipeline comprises five main processing stages:

\begin{enumerate}
  \item \textbf{Data Acquisition}: Automated retrieval of seismic waveforms 
  from data servers. To prevent server timeouts, downloads were segmented into 
  days of interval.

  \item \textbf{Phase Picking}: Machine learning-based detection of P and S 
  wave arrivals using the SeisBench framework. This stage leveraged GPU 
  acceleration for neural network inference.

  \item \textbf{Phase Association}: Grouping of detected phases into coherent 
  seismic events using two complementary algorithms:
  \begin{itemize}
    \item GaMMA: A probabilistic Gaussian Mixture Model associator
    \item PyOcto: An octree-based, velocity-model-aware association algorithm 
    optimized for computational efficiency
  \end{itemize}

  \item \textbf{Event Location}: Non-linear hypocenter determination using the 
  NonLinLoc software package, employing probabilistic global search algorithms 
  with a 1D layered velocity model.

  \item \textbf{Magnitude Estimation}: Computation of local magnitudes ($M_L$) 
  from Wood-Anderson simulated waveforms using regional attenuation
  corrections.
\end{enumerate}

\subsection{Workflow Automation}
The Makefile-based automation system provides the following capabilities:

\begin{itemize}
  \item \textbf{Modular Execution}: Individual pipeline stages can be executed 
  independently, enabling selective reprocessing

  \item \textbf{Configuration Management}: The pipeline uses Hydra 
  (\cite{Hydra}), an open-source framework developed by Meta for composable 
  configuration management. Hydra organizes all processing parameters, such as 
  detection thresholds, model selections, station lists, and velocity-model 
  paths, into hierarchical YAML files that can be overridden from the command 
  line at runtime without modifying source code. This enables systematic 
  parameter sweeps (\textit{e.g.} varying detection thresholds across 0.1, 0.2, 
  0.3) and ensures that every pipeline run is fully reproducible from its 
  configuration alone.

  \item \textbf{Disk Space Optimization}: Symbolic links are used to share 
  intermediate results between processing stages, reducing storage requirements

  \item \textbf{Multi-Year Processing}: Parameterized targets support 
  processing of multiple years with consistent configurations
\end{itemize}

\section{Parallelization Strategy}

The deep learning pipeline is parallelized and optimized for performance on the 
Leonardo cluster using multiple parallelization paradigms.

\subsection{Parallel Execution Framework}
The pipeline employs a hybrid parallelization approach:

\begin{itemize}
  \item \textbf{CUDA}: GPU acceleration for neural network inference during the 
  phase picking stage

  \item \textbf{OpenMP}: Thread-level parallelism for shared-memory operations 
  within individual nodes
\end{itemize}

\subsection{Thread Optimization}
The resource allocation strategy ensures full utilization of available CPU 
cores by maintaining the constraint:

\begin{equation}
  N_{\text{tasks}} \times N_{\text{threads}} = N_{\text{cores}}
\end{equation}

where $N_{\text{tasks}}$ represents the number of MPI processes (or GPU 
instances), $N_{\text{threads}}$ represents the number of OpenMP threads per 
task, and $N_{\text{cores}}$ is the number of available CPU cores per node (32 
on the Leonardo booster partition). This approach maximizes CPU utilization 
while allowing flexibility in the parallelization strategy.

\subsection{Execution Modes}
The system supports two execution modes selected automatically based on the 
task count:

\begin{itemize}
  \item \textbf{Serial Mode}: For single-task jobs, direct Python execution 
  without MPI overhead

  \item \textbf{Parallel Mode}: For multi-task jobs, MPI-based parallel 
  execution using \texttt{mpirun}
\end{itemize}

\section{Parallel I/O}

The input/output operations are parallelized and optimized for performance on 
the Leonardo cluster.

\subsection{I/O Parallelization}
The I/O operations leverage:

\begin{itemize}
  \item \textbf{Lustre File System}: High-speed parallel storage optimized for 
  large-scale scientific computing

  \item \textbf{Striping Configuration}: Data files are striped across multiple 
  Object Storage Targets (OSTs) for improved throughput
\end{itemize}

\subsection{Task Distribution}
The task distribution strategy is optimized to:

\begin{itemize}
  \item Minimize data transfer between nodes
  \item Balance computational load across nodes
  \item Partition seismic waveforms into temporal chunks for parallel 
  processing
  \item Ensure even distribution of workload across all available resources
\end{itemize}

\section{Job Submission Framework}

A template-based job submission wrapper (\texttt{Analyze.sh}) was developed to 
streamline the submission process on the SLURM scheduler.

\subsection{Submission Workflow}
The wrapper performs the following operations:

\begin{enumerate}
  \item Modifies a template SLURM script with the required resource 
  specifications (nodes, tasks, CPUs, GPUs)

  \item Configures environment variables for thread management 
  (\texttt{OMP\_NUM\_THREADS}, \texttt{NUMBA\_NUM\_THREADS})

  \item Submits the configured job to the SLURM scheduler

  \item Restores the template to its original state for subsequent submissions
\end{enumerate}

\subsection{Dynamic Resource Configuration}
This approach enables rapid reconfiguration of computational resources without 
manual editing of job scripts, facilitating experimentation with different 
parallelization strategies.

\section{Scalability Considerations}

The pipeline was designed with scalability as a primary objective.

\subsection{Design Principles}
Key scalability features include:

\begin{itemize}
  \item \textbf{Temporal Segmentation}: Data downloads are segmented by time to 
  prevent server timeouts and enable parallel acquisition

  \item \textbf{GPU Acceleration}: Machine learning inference on GPUs reduces 
  phase picking time by orders of magnitude compared to CPU-only execution

  \item \textbf{Modular Architecture}: Independent tasks can be distributed 
  across multiple nodes

  \item \textbf{Shared Storage}: Output catalogs are stored on shared 
  filesystems accessible to all compute nodes
\end{itemize}

\subsection{Resource Allocations}
Table~\ref{tab:hpc_resources} summarizes the typical resource allocations used 
for each pipeline stage.

\begin{table}[h]
\centering
\caption{Typical HPC resource allocations for pipeline stages}
\label{tab:hpc_resources}
\begin{tabular}{lccc}
\hline
\textbf{Pipeline Stage} & \textbf{Nodes} & \textbf{Tasks/Node} & \textbf{Threads/Task} \\
\hline
Data Acquisition & 1 & 1 & 32 \\
Phase Picking (SeisBench) & 1 & 1 & 32 \\
Phase Association (GaMMA) & 1 & 1 & 32 \\
Phase Association (PyOcto) & 1 & 1 & 32 \\
Event Location (NonLinLoc) & 1 & 1 & 32 \\
\hline
\end{tabular}
\end{table}

% =============================================================================
% End of HPC Implementation Chapter
% =============================================================================